%------------------------------------------------------------------------------%
\begin{question}
\end{question}

%------------------------------------------------------------------------------%
\begin{resolution}{Solução}
\end{resolution}

%------------------------------------------------------------------------------%
\endinput

% ============================================================
\begin{frame}{Exemplo 4}
Considere
\[
A=(2,-1,4)
\]
e
\[
B=(-1,3,2)
\]

Calcule
\[
\norm{B-A}
\]
e interprete geometricamente o resultado.
\end{frame}

% ============================================================
\begin{frame}{Exemplo 4 --- Solução}
\begin{align*}
B-A
&=
(-1,3,2)-(2,-1,4)
\\
&=
(-3,4,-2)
\end{align*}

Logo,
\begin{align*}
\norm{B-A}
&=
\sqrt{(-3)^{2}+4^{2}+(-2)^{2}}
\\
&=
\sqrt{9+16+4}
\\
&=
\sqrt{29}
\end{align*}

Portanto,
\[
d(A,B)=\sqrt{29}
\]

Geometricamente, $\sqrt{29}$ é o comprimento do segmento que liga $A$ a $B$.
\end{frame}
